\subsection{Future Work}
\label{sec-future-work}

Our future work 
can be divided into three thrusts consisting
of seven {\em initial} tasks:
\begin{enumerate}
\item {\em Learning method improvement thrust} (Tasks $A$, $B$, and $C$).

We will improve the effectiveness of
$\MEGA$ and $\SVMActive$, and possibly lead to
the development of new online learning algorithms
(e.g., Bayes Point Machines with Active Learning). We plan to 
investigate algorithms to perform {\em adaptive sampling}, and
to explore {\em co-training} using $\MEGA$ and $\SVMActive$,
and study integrating active learning with
Bayes Point Machines. 

\item {\em Chen's thrust} (Tasks $D$, $E$, $F$).

% $Id: chen.tex,v 1.2 2001/10/29 17:38:27 chenli Exp $

\subsection{Heterogeneous data integration}

\subsubsection*{Task 1: Extracting semi-structured data from web pages}

To support cross-modality search, we need to extract structured
information from web pages.  This extraction is also critical to many
other applications.  For instance, comparison-shopping engines need to
retrieve product information from their underlying sources.  In addition,
extracting web information effectively can meet the growing demand in
making web content available to mobile devices (e.g., PDAs and cell
phones), which have limited display and network bandwidth.

Many approaches have been proposed on extracting information from web
pages.  Many applications need software components called {\em scrapers}
to do the extracting automatically.  We see that the data is going through
a full cycle: data is stored in structured databases at publisher sites,
published in unstructured formats, and scraped and stored in structured
databases again at the consumer end.  Scrapers are typically built using
rule-based techniques.  Often, these techniques involve a human spending a
lot of time on the sites, figuring out the format and layout of the pages,
and then writing custom rules for scraping the data.  Apart from being
time consuming and error prone, these techniques are not robust to changes
in the webpage structure.

In our study, we plan to tackle this problem by analyzing the regularities
shared by different pages of the same web site.  For instance,
Figure~\ref{fig:bookpages} shows two book pages at Amazon.com.  The figure
clearly indicates the similarities of the two pages.  They share the same
layout, in which some portions have same contents.  One important
observation is that some interesting information, such as book title,
authors, price, and availability, are included in the differences.

\begin{figure}[htfb]
\begin{center}
\begin{tabular}[htfb]{ll}
\epsfxsize=200pt
\epsfysize=150pt
\subfigure[Book 1.]{\epsfbox{amazon1.eps}}
\epsfxsize=200pt
\epsfysize=150pt
\subfigure[Book 2.]{\epsfbox{amazon2.eps}}
\end{tabular}
\end{center}
\caption{Two book pages at Amazon.com.}
\label{fig:bookpages}
\end{figure}

In general, a lot of Web information is stored in the form of highly
structured, often dynamically generated pages.  These pages are typically
generated by plugging in data from a structured backend source (relational
database) into HTML templates.  Our approach is to infer the structure of
these HTML templates used to generate the webpages by examining many
webpages.  By analyzing these templates we can extract useful information
effectively.  Our approach works in the following steps:

\begin{enumerate}
  \item Choose some representative pages from a web site that looks
  similar.

  \item Extract the common template shared by these pages.

  \item Show the differences to users.  Let users choose the values that
  are useful.

  \item Modify the template.  Apply the new template to other pages to
  extract useful information to populate the result database.
\end{enumerate}

Our approach will work in the case where more than one template is used to
generate the pages.

\subsubsection*{Task 2: Constructing Meta-data for Multimedia Objects}

It is known that content-based search on multimedia information (e.g.,
images, audio, and video) has limitations, since it is hard to extract
features that can match concepts understandable by human being.  Hence
constructing meta data for multimedia objects from web pages is very
important to support effective cross-modality search.  For instance, given
an image of a painting of Monet, it is hard, if not impossible, for a
content-based system to extract the information such as ``this painting
was drawn by Monet,'' and ``it is an art of impressionism.''  However, by
analyzing the surrounding text and anchor text on links to this page, we
are able to compute this meta data.

We study how to construct the meta data of a multimedia object effectively
by using its surrounding text and anchor text of its links.  The following
are some challenges and approaches we will take:

(1) Extracting representative keywords in surrounding text.  In what range
    should we extract text that might provide representative keywords for
    a given object?  How to extract these keywords?  We need to compute a
    ``feature'' of the surrounding texts of an object, such as the range
    of the text, the sequence of each word, etc.  We need to choose the
    feature carefully, and then choose the important keywords
    correspondingly.

(2) Computing meta data using common patterns.  Similar to Task 1, we can
    use the fact that the ways multimedia objects are placed in web pages
    share some common layouts.  For instance, often people put a caption
    below an image to describe the content of the image.  If we can
    extract the common patterns these objects are placed, we can use these
    patterns to help us extract useful meta data for other objects.  That
    is, we start from some seed objects, together with their meta data,
    and see how they appear in web pages.  Then we compute the patterns of
    these objects, and apply these patterns to web pages to extract more
    objects with their meta data.  We iteratively extract patterns and
    objects (with their meta data).

    This approach is similar to the idea in \cite{Brin}.  However, the
    following challenges make our problem unique:

\begin{itemize}
   \item How to represent the layout of an object in a page?  To what
   granularity do we want to describe the layout?  This problem also
   exists in Task 1.  In \cite{Brin}, we can use a regular expression to
   represent an entity (e.g., a book title) is mentioned in a sentence.
   Here we are dealing with multimedia objects at a higher level.

   \item How to identify a multimedia object consistently in different
   pages?  For example, the same image could be in different formats in
   different pages, such as JPG, GIF, and BMP.  Some images could have
   transformed slightly.  Our approach should be able to identify the
   multiple instances of the same object in different pages, even if they
   have different binary codes.
\end{itemize}

(3) Mining anchor text on links to construct meta data.  Anchor text on
    links to a multimedia object provides very rich content about the
    object.  In addition, different web pages have different ``authority''
    in identifying the content of an object.  So we can utility the
    results of Hub/Authority or PageRank to weight the importance of the
    anchor text from different pages.  By using a weighting function on
    different links to an object, we can compute the meta data accurately.

\subsubsection*{Task 3: Computing cross-modality result using ranking functions}

To do a cross-modality query, we often need to combine information of text
and multimedia objects.  Thus it is critical to have an effective ranking
function for these different types of information.  For instance, suppose
a user wants to find all the web pages that contain a Monet painting image
and are related to impressionism.  To answer this query, we need to access
text information to find pages related to impressionism, and access image
information to find images of Monet paintings.  We should choose the right
weighting function for these different types of information to compute the
right relevance of each page to a query.

In addition, given a ranking function, how to compute the results very
efficiently is also a thrust of our research.  One challenge is that the
data size is large.  Currently there are more than 1.6 billion pages on
yhr Web.  Furthermore, the number of pages is growing very rapidly.  Our
searching algorithm should be efficient and scalable.  In the case where
it takes too much time to compute the answer, the algorithm will return
some approximate answers without compromising the accuracy too much.

The algorithm can use user feedback to improve its efficiency and
accuracy.  The user could either explicitly specify what pages are good or
bad, or click on pages, or the system will automatically detect what pages
are more preferred by the user.  Then the algorithm will adjust its
weighting function based on the user feedback to improve the result
quality for future queries.  In this sense, the algorithm is {\\em
personalized}, since it automatically adjusts its weighting function based
on the feedback from different users.

In addition, to support the weighting function, we need to design
efficient indexing structures for different modalities.

\subsubsection*{Related work}

{\em Data Integration}: The rapid growth of the Internet is giving us
access to an unprecedented number of heterogeneous information sources.
Many information-integration systems (e.g., TSIMMIS \cite{Chawathe94-1},
the Information Manifold \cite{Levy96-1}, Garlic \cite{Roth97-1},
Infomaster \cite{Genesereth97-1}, Disco \cite{Tomasic98-1}, Tukwila
\cite{Ives99-1}, and InfoSleuth \cite{Bayardo97-1}) have been proposed to
support seamless access to these heterogeneous data sources.  These
systems build {\em wrappers} on web sources to extract information from
web pages.

{\em Multimedia Ranking}: Several techniques (\cite{Fagin,Gravano,Chang})
have been proposed on how to rank information from different sources.


%Related work: web-page extraction.  Wrapper generation.  Chen's work on
%integration.


\item {\em Junghoo's thrust} (Tasks $G$, $H$, and $I$).
\end{enumerate}

\subsubsection*{\underline{Task $A$: Adaptive Sampling}}

Both $\MEGA$ and $\SVMActive$ rely on effective sampling 
to select the most informative training data to query endusers.
The sampling methods of these schemes can be
improved in the following two directions:

First, both schemes use the same sampling method from one
user iteration to another.  Our new hypothesis is that
in different stages of the learning process, the learner
may want to place more weight on {\em exploring} or
on {\em exploiting}.
Intuitively, when the learner is used to explore the
uncertainty space, a more aggressive sampling approach should be used. 
When the learner needs to perform fine-grained concept refinement, 
it should employ a more
conservative sampling scheme such as the traditional ``nearest neighbor.''
One possible criterion for determining whether to
employ a more
aggressive or more conservative sampling scheme 
can be the number of instances in the
positive-labeled pool.  When 
few samples have been labeled positive, we would 
want the learner 
to be more explorative in the uncertain feature space.
For example, if a red flower is
the only positive-labeled instance, this does not necessarily
mean that the user's query concept is red flowers.  It is
possible that the query concept is ``something red'' or
``flowers of other kinds.''  In this case,
the sampling scheme 
should be more explorative.  Conversely, if the learner has
collected a fairly large number of positive-labeled instances
that are similar to each other, then the preferred sampling
scheme would be 
fine-grained refinement.  We will first conduct extensive
experiments to examine whether the above hypothesis
holds.  If it does hold,
we will develop a theoretical framework for adaptively
determining the appropriate sampling scheme.

The second research
direction is to design more effective sampling schemes.
Once a user determines that the learner 
should be more explorative, 
it needs a robust sampling method to find the most appropriate samples
to quickly reduce the uncertain space.
$\SVMActive$ currently employs Tong \&
Koller's {\em simple method} \cite{TK00-b}, which selects all samples 
nearest the hyperplane that divides the positive-
and negative-labeled instances.  However, this theorem
has been found to work optimally only when
samples are taken one at a time.  But it is
necessary to present a reasonable number
of samples to query users in each user iteration, since
most users will not engage in even a moderate
number of feedback rounds.  

One of our conjectures is considering {\em transduction}.
{Transduction} takes advantage of not only the labeled instances but
also the unlabeled instances to learn a user's query concept.
Intuitively, the unlabeled instances are likely to be clustered. This
tendency provides us with some extra information which we may use to
refine our idea of the user's query concept.  Transduction
has been successfully applied to regular passive SVMs for text
classification~\cite{Joachims:ICML99} showing that using unlabeled
data can indeed improve performance. An exciting 
direction for new research is to see
whether transduction can be applied to the image retrieval domain and
furthermore, whether it can also be combined with active learning.

Our preliminary experiments with transduction agrees with 
some observations \cite{GHK00} that ``transduction consistently behaves 
inconsistently.''  However, with active learning, 
our results show that transduction can be made more effective
through {\em confirmation}. 
We can present nearest neighbors of positive/negative samples
to the user for confirmation.  
Depending on user feedback, our algorithm can adjust
how aggressively it conducts transduction.

\subsubsection*{\underline{Task $B$: Co-training}}

$\MEGA$ is an inductive-based active learning algorithm, and $\SVMActive$
is a margin-based one.
$\MEGA$ and $\SVMActive$ can complement each other in 
at least two ways.  
First, $\SVMActive$ requires seeding at least one
positive and one negative sample to initiate learning
a query concept.  Finding a positive sample can be very difficult 
when the concept complexity is high.  For instance, 
finding a flower image in a large and diversified 
dataset containing thousands of image categories
can take quite a while.  $\MEGA$, on the other hand, does not
require seeding.  If all of the images
generated in the first round of feedback are found irrelevant,
it will use the negative images to reduce the set of potentially
relevant images substantially; thus the
probability is higher that a relevant image will be
generated in the next round.
The logical design choice is using $\MEGA$ to bootstrap
a learning session until at least one sample
is labeled positive.

After identifying some positive and some negative labeled instances,
$\SVMActive$ employs a well-founded 
sample selection method (at least so far) for selecting 
data objects nearest the hyperplane 
that divides positive from
negative instances. $\MEGA$, on the other hand, 
depends more on heuristics in choosing samples
that can eliminate the maximum expected number of
disjunctive terms.  
$\MEGA$ and $\SVMActive$ present two distinct views
of each example, and this fits the co-training model
proposed by Blum and Mitchell \cite{BM98}.
It is logical for two schemes
to exchange sampling information, and if these
methods can be adapted to work
in a complementary manner, even more effective
learning may result.

\subsubsection*{\underline{Task $C$: Active Learning with Bayes Point Machines}}

\begin{figure}[h]
\centerline{
\begin{tabular}{ccc}
\psfig{figure=vspace1.eps,height=1.5in,angle=270} &
\hspace{-0.1in}
\psfig{figure=vspace2.eps,height=1.5in,angle=270} \\
(a) & (b)
\end{tabular}}
\footnotesize{
\caption{(a) Version space duality. The surface of the hypersphere
represents unit weight vectors. Each of the two hyperplanes
corresponds to a labeled training instance. Each hyperplane restricts
the area on the hypersphere in which consistent hypotheses can
be situated. Here version space is the surface segment of the hypersphere
closest to the camera. (b) An SVM classifier in version space. The
dark embedded sphere is the largest radius sphere whose center lies in
version space and whose surface does not intersect with the
hyperplanes. The center of the embedded sphere corresponds to the SVM,
its radius is the margin of the SVM in $\setF$, and the training points
corresponding to the hyperplanes that it touches are the support
vectors.}
}
\label{fig:duality}
\label{fig:simple1}
\end{figure}

From a Bayesian perspective Support Vector Machines choose the 
hypothesis corresponding to the largest possible hypersphere 
that can be inscribed in version space, i.e. in the 
space of all consistent hypotheses given a training set \cite{SVM01}.
Given a set of labeled training data and a Mercer kernel $K$, there is
a set of hyperplanes that separate the data in the induced feature
space $\setF$. We call this set of consistent hyperplanes or {\em
hypotheses} the {\em version space}~\cite{Mitchell:AI82}. In other
words, hypothesis $f$ is in version space if for every training
instance $\vc{x}_i$ with label $y_i$ we have that $f(\vc{x}_i) > 0$ if
$y_i = 1$ and $f(\vc{x}_i) < 0$ if $y_i = -1$.
Clearly, by definition, points in $\setW$ correspond to hyperplanes in
$\setF$. The intuition behind the converse is that observing a
training instance $\vc{x}_i$ in feature space restricts the set of
separating hyperplanes to ones that classify $\vc{x}_i$ correctly. In
fact, we can show that the set of allowable points $\vc{w}$ in $\setW$
is restricted to lie on one side of a hyperplane in $\setW$.  More
formally, to show that points in $\setF$ correspond to hyperplanes in
$\setW$, suppose we are given a new training instance $\vc{x}_i$ with
label $y_i$. Then any separating hyperplane must satisfy
$y_i(\innerprod{\vc{w}}{\Phi(\vc{x}_i)}) > 0$. Now, instead of viewing
$\vc{w}$ as the normal vector of a hyperplane in $\setF$, think of
$y_i \Phi(\vc{x}_i)$ as being the normal vector of a hyperplane in
$\setW$. Thus $y_i(\innerprod{\vc{w}}{\Phi(\vc{x}_i)}) =
\innerprod{\vc{w}}{y_i \Phi(\vc{x}_i}) > 0$ defines a half-space in
$\setW$.  Furthermore $\innerprod{\vc{w}}{y_i \Phi(\vc{x}_i}) = 0$
defines a hyperplane in $\setW$ that acts as one of the boundaries to
version space $\setV$. Notice that version space is a connected region
on the surface of a hypersphere in parameter space. 

Figure~\ref{fig:duality}(a) presents an example.
SVMs find the center of the largest radius hypersphere whose center
can be placed in version space and whose surface does not intersect
with the hyperplanes corresponding to the labeled instances, as in
Figure~\ref{fig:duality}(b). It can be easily shown that the
hyperplanes touched by the maximal radius hypersphere
correspond to the support vectors and that the radius of the
hypersphere is the margin of the SVM.

An alternative and potentially better approach is to construct the 
hypothesis using the whole of version space. This is achieved by using a 
Bayes Point Machine that finds the midpoint of the region of 
intersection of all 
hyperplanes bisecting version space into two halves of equal volume (the 
Bayes point) \cite{Herbrich+al,HG01,HGC01}. It is known that the centre of 
mass of version space 
approximates the Bayes point. The initial work of
Herbrich and Graepel estimates the centre of 
mass by averaging over the trajectory of a billiard 
ball bouncing in version space, sampling version space using a 
Gibbs sampler and approximative methods using the perceptron 
learning algorithm \cite{GH01}.
This method, however, is not efficient.
We plan to investigate new algorithms that can find the 
centre of mass more efficiently.  We also plan to
implement active learning with Bayes Point Machines and
compare its performance with $\SVMActive$.

